\documentclass[11pt]{article}

\usepackage[T1]{fontenc}
\usepackage[utf8]{inputenc}
\usepackage{lmodern}
\usepackage[margin=1in]{geometry}
\usepackage{amsmath,amssymb,amsthm,mathtools}
\usepackage[hidelinks]{hyperref}

\newtheorem{theorem}{Theorem}
\newtheorem{proposition}{Proposition}
\newtheorem{corollary}{Corollary}
\theoremstyle{remark}
\newtheorem{remark}{Remark}
\newtheorem{example}{Example}

\newcommand{\rank}{\operatorname{rank}}
\newcommand{\RePart}{\operatorname{Re}}
\newcommand{\dd}{\mathrm{d}}

\title{A Rank and Signature Obstruction to Lorentzian Metric Emergence\\
from Positive-Sign Scalar Gradients}
\author{Richard J. Reyes}
\date{October 26, 2026}

\begin{document}
\maketitle

\begin{abstract}
We establish a pointwise algebraic obstruction to constructing a four-dimensional
Lorentzian metric from the real part of the derivative product of a single
complex scalar field.  Writing the scalar as two real components shows that
its proposed symmetric tensor is a Gram form with rank at most two and is
positive semidefinite.  Consequently, the tensor is degenerate and cannot
have Lorentzian signature.  More generally, any finite sum of real
scalar-gradient outer products with nonnegative weights is positive
semidefinite, regardless of how many gradient directions are supplied.
The conclusion is specific to this positive-sign gradient construction;
it does not exclude composite or effective metrics involving different
tensorial ingredients or independent geometric structure.
\end{abstract}

\section{Setup and main theorem}

Let $M$ be a smooth real four-dimensional manifold and let
$\psi\in C^1(M;\mathbb C)$ be a complex scalar field.  Write
\begin{equation}
  \psi=a+ib, \qquad a,b\in C^1(M;\mathbb R).
\end{equation}
Define the real symmetric covariant $(0,2)$-tensor
\begin{equation}\label{eq:defT}
  T_{\mu\nu}
   :=\RePart\bigl(\partial_\mu\overline{\psi}\,
                         \partial_\nu\psi\bigr).
\end{equation}
This definition is coordinate independent because derivatives of scalar
fields are covectors.

\begin{theorem}[Single-complex-scalar rank and signature obstruction]
\label{thm:single}
At each point $p\in M$, the bilinear form $T_p$ defined in
\eqref{eq:defT} satisfies
\begin{align}
  \rank T_p &\leq 2, \label{eq:rankbound}\\
  T_p(V,V)&\geq 0 \quad\text{for every }V\in T_pM, \label{eq:psd}\\
  \dim\ker T_p&\geq 2. \label{eq:kernelbound}
\end{align}
Here $\ker T_p=\{V\in T_pM:T_p(V,W)=0\text{ for all }W\in T_pM\}$
is the radical of the bilinear form.  In particular, $\det T_p=0$
in any coordinate basis, and $T_p$ cannot be a Lorentzian metric.
\end{theorem}

\begin{proof}
Since $\overline{\psi}=a-ib$,
\begin{align}
 \partial_\mu\overline{\psi}\,\partial_\nu\psi
 &= (\partial_\mu a-i\partial_\mu b)
    (\partial_\nu a+i\partial_\nu b)\\
 &= \partial_\mu a\,\partial_\nu a
  + \partial_\mu b\,\partial_\nu b
  + i(\partial_\mu a\,\partial_\nu b
      -\partial_\mu b\,\partial_\nu a).
\end{align}
Taking the real part gives the exact identity
\begin{equation}\label{eq:sumouter}
  T_{\mu\nu}=\partial_\mu a\,\partial_\nu a
                +\partial_\mu b\,\partial_\nu b.
\end{equation}
Fix $p\in M$ and a coordinate basis for $T_pM$.  Set
$A_\mu=(\partial_\mu a)_p$ and $B_\mu=(\partial_\mu b)_p$.
Introduce the real $4\times 2$ matrix
\begin{equation}
  J=\begin{pmatrix}
   A_0&B_0\\ A_1&B_1\\ A_2&B_2\\ A_3&B_3
  \end{pmatrix}.
\end{equation}
Equation \eqref{eq:sumouter} is the Gram factorization
\begin{equation}\label{eq:gram}
  [T_p]=JJ^{\mathsf T}.
\end{equation}
It follows that
\begin{equation}
 \rank T_p=\rank(JJ^{\mathsf T})\leq \rank J\leq 2,
\end{equation}
proving \eqref{eq:rankbound}.  In particular, all $3\times3$
minors of $[T_p]$ vanish, and its $4\times4$ determinant vanishes.

For any real tangent vector $V=V^\mu\partial_\mu|_p$,
\begin{align}\label{eq:squares}
 T_p(V,V)
 &= (V^\mu\partial_\mu a)^2
   +(V^\mu\partial_\mu b)^2\\
 &=\|J^{\mathsf T}[V]\|_{\mathbb R^2}^2\geq 0.
\end{align}
This proves positive semidefiniteness.  Finally, rank--nullity
gives
\begin{equation}
 \dim\ker T_p=4-\rank T_p\geq 2.
\end{equation}
The kernel consists precisely of the tangent vectors annihilated by
both $\dd a|_p$ and $\dd b|_p$, as follows from
\eqref{eq:squares} and \eqref{eq:sumouter}.

A Lorentzian metric in four dimensions is a nondegenerate symmetric
bilinear form of inertia $(1,3)$ or $(3,1)$: it is indefinite and
has rank four.  By contrast, $T_p$ is both degenerate and positive
semidefinite.  Either property alone excludes Lorentzian signature.
These properties are invariant under changes of coordinates
(the inertia statement follows from Sylvester's law of inertia).
Therefore $T_p$ cannot be a Lorentzian metric.
\end{proof}

\section{Positive-sign multiplets}

\begin{proposition}[Nonnegative scalar-gradient sums]\label{prop:multiplet}
Let $\phi^1,\dots,\phi^N\in C^1(M;\mathbb R)$ and let
$w_A:M\to[0,\infty)$ be nonnegative weights.  Define
\begin{equation}\label{eq:multi}
 S_{\mu\nu}:=\sum_{A=1}^N w_A\,
   \partial_\mu\phi^A\,\partial_\nu\phi^A.
\end{equation}
At every $p\in M$,
\begin{equation}
 \rank S_p\leq\min\{N,4\},\qquad
 S_p(V,V)\geq 0\quad\text{for all }V\in T_pM.
\end{equation}
Thus no such tensor is Lorentzian, for any finite $N$.
\end{proposition}

\begin{proof}
At a fixed $p$, define the real $4\times N$ matrix
\begin{equation}
 K_{\mu A}:=\sqrt{w_A(p)}\,(\partial_\mu\phi^A)_p.
\end{equation}
Then $[S_p]=KK^{\mathsf T}$, and hence
\begin{equation}
 \rank S_p\leq\rank K\leq\min\{N,4\}.
\end{equation}
Moreover, for every real $V\in T_pM$,
\begin{equation}
 S_p(V,V)=\sum_{A=1}^N w_A(p)
 \bigl(V^\mu\partial_\mu\phi^A\bigr)^2\geq 0.
\end{equation}
A Lorentzian form is indefinite; therefore $S_p$ cannot be
Lorentzian, even if it has rank four.
\end{proof}

\begin{corollary}[Rank budget is necessary but not sufficient]
Within the construction \eqref{eq:multi}, at least four independent
real gradient directions are necessary to obtain a rank-four tensor
at a point.  Having four or more such directions cannot by itself
produce Lorentzian signature when the weights remain nonnegative.
\end{corollary}

\section{Sharp examples and the role of additional structure}

\begin{example}[Sharpness of the rank-two bound]
In coordinates $(x^0,x^1,x^2,x^3)$ on $\mathbb R^4$, choose
$a=x^1$ and $b=x^2$.  Then
\begin{equation}
 [T]=\operatorname{diag}(0,1,1,0),
\end{equation}
so $\rank T=2$ exactly.
\end{example}

\begin{example}[Four gradients give full rank, not Lorentzian signature]
Take $\phi^A=x^{A-1}$ for $A=1,2,3,4$ and $w_A=1$.  Then
\begin{equation}
 [S]=\operatorname{diag}(1,1,1,1),
\end{equation}
which is full rank but positive definite, not Lorentzian.
\end{example}

\begin{example}[Indefinite internal bilinear form]
Suppose instead one introduces an internal tensor
$\eta_{AB}=\operatorname{diag}(-1,1,1,1)$ and defines
\begin{equation}
 G_{\mu\nu}:=\eta_{AB}\,
  \partial_\mu\phi^A\,\partial_\nu\phi^B.
\end{equation}
For $\phi^A=x^A$ with indices $A=0,1,2,3$, one gets
\begin{equation}
 [G]=\operatorname{diag}(-1,1,1,1).
\end{equation}
Thus Lorentzian signature is algebraically possible once an
indefinite internal bilinear form is supplied.  The internal form
is additional input, not a consequence of \eqref{eq:defT}.
\end{example}

\begin{remark}[Scope of the no-go statement]
The theorem excludes identifying the specific tensor
$\RePart(\dd\overline\psi\otimes\dd\psi)$, or any nonnegative
scalar-gradient sum of the form \eqref{eq:multi}, with a Lorentzian
metric.  It does not prove that no theory involving a complex scalar
can possess an effective Lorentzian metric.  In particular, it does
not address constructions using an independent metric, indefinite
contributions, disformal combinations, additional tensor fields,
or other geometric data.  Nor does it establish Einstein dynamics,
diffeomorphism gauge structure, a massless spin-two mode, or a
microscopic theory of gravity.
\end{remark}

\section{Conclusion}
The proposed tensor is a Gram form of only two real covectors:
\begin{equation}\boxed{
 T_{\mu\nu}=\partial_\mu a\,\partial_\nu a
             +\partial_\mu b\,\partial_\nu b,
 \quad\rank T\leq2,\quad T(V,V)\geq0.
}\end{equation}
Thus it is everywhere degenerate and positive semidefinite, and
cannot be a four-dimensional Lorentzian metric.  More positive-sign
scalar gradients can remove the rank deficiency, but cannot change
the positive-semidefinite signature.  Lorentzian metric emergence
therefore requires a different tensor construction, an indefinite
contribution, or additional geometric structure.

\paragraph{Research context.}
This observation is applied to the tensor/tetrad emergence question
in the foundational audit of the geometry-of-resonance program:
\url{https://github.com/rickyjreyes/geometry_of_resonance/blob/main/docs/FOUNDATIONAL_STAGE2_2026-10-14.md}.
The underlying rank and Gram-matrix facts are standard linear algebra;
the no-go conclusion concerns the specified candidate construction.

\end{document}
